% =====================================================================
% Thesis section: High-Frequency Adversarial Supervision for SAR-to-Optical
% Translation under Geometric Misregistration (LLW-Former v0.5.x / llwt_v5)
%
% Drop-in chapter/section. Requires: amsmath, booktabs, graphicx, siunitx (optional).
% Figure/checkpoint placeholders marked <FILL>.
% =====================================================================

\section{High-Frequency Adversarial Supervision under Geometric Misregistration}
\label{sec:hfd}

\subsection{Motivation and problem analysis}
\label{sec:hfd:motivation}

The LLW-Former generator (Section~\ref{sec:llwformer}) attains competitive low-frequency
fidelity on Sentinel-1/2 (SEN1-2) SAR$\rightarrow$optical translation, yet its outputs remain
perceptibly soft. To locate the bottleneck precisely we profiled a trained checkpoint on a
held-out set of \num{320} aligned validation patches and decomposed the error by spatial
frequency and by Haar wavelet sub-band (Table~\ref{tab:errorprofile}).

\begin{table}[t]
\centering
\caption{Error profile of the baseline generator on held-out aligned validation patches.
Relative band error is the per-band $L_1$ normalised by that band's energy; the radial power
ratio is $P_\text{fake}/P_\text{opt}$ in each spatial-frequency octave (target $=1.0$).}
\label{tab:errorprofile}
\begin{tabular}{lcc}
\toprule
Quantity & Value & Interpretation \\
\midrule
Luminance / chrominance error & $9.43$ / $4.5$ & colour is \emph{not} the bottleneck \\
Global colour bias & $-0.003$ & no cast \\
Rel.\ error LL (low-freq.) & $0.33$ & coarse structure recovered \\
Rel.\ error detail (LH,HL,HH) mean & $\geq 1.3$ & detail phase-\emph{decorrelated} \\
Detail / LL relative error ratio & $4.2$ & detail $4\times$ worse than LL \\
Mid-freq.\ power ratio & $0.66$ & \textbf{blur} (missing mid structure) \\
High-freq.\ power ratio & $1.14$ & incoherent over-energy (noise) \\
\bottomrule
\end{tabular}
\end{table}

Three facts emerge. (i) The deficit is \emph{structural, not chromatic}: chrominance error is
roughly half the luminance error and the global colour bias is negligible. (ii) The model
recovers the low-frequency band (LL relative error $0.33$) but its detail sub-bands are
phase-decorrelated (relative error $>1$, i.e.\ the error exceeds the band's own energy). (iii)
Spectrally, the generator \emph{under-produces} mid frequencies (power ratio $0.66$, the blur
signature) while \emph{over-producing} incoherent high frequencies (ratio $1.14$). The worst
patches are low-texture water/river scenes, consistent with the one-to-many nature of the
mapping: from speckled SAR, the exact optical micro-texture of a smooth surface is
underdetermined.

A second, compounding factor is geometric: paired SEN1-2 patches are misregistered. Estimating
per-pair rigid alignment (gradient-domain phase correlation refined by ECC) over \num{300}
sampled pairs of the training scenes yields a median shift of \SI{1.23}{px} (p90
\SI{5.3}{px}; $58\%$ of pairs exceed \SI{1}{px}, $31\%$ exceed \SI{2}{px}), plus a small
rotational component (mean $0.38^\circ$). Such misregistration makes position-exact
supervision of high-frequency content unreliable.

\subsection{Methodology: hypothesis-driven elimination}
\label{sec:hfd:method}

Rather than commit compute to a full \SI{200}{epoch} run per idea, each candidate was first
validated (or refuted) by a targeted, inference- or short-optimisation-level experiment. This
discipline produced two negative results that are contributions in their own right and that
directly motivate the final method.

\subsubsection{Negative result 1: dense learned registration is infeasible for SAR$\to$optical}
\label{sec:hfd:neg1}

The natural response to misregistration is to learn a dense deformation field $\phi$ that warps
the target into the generator's geometry before computing fidelity losses (a self-aligning
translator). We implemented such an aligner, estimating $\phi$ from the speckle-robust LL Haar
band, and found it \emph{inert}: across \num{130} training epochs $\lvert\phi\rvert$ remained
pinned at ${\sim}\SI{0.07}{px}$ (the bf16 sampling-noise floor), confirmed identical in fp32.
Controlled probes isolated the cause. With the regulariser removed and the generator frozen,
the fidelity gradient could drive $\phi$ only to ${\sim}\SI{0.2}{px}$, far below the measured
\SI{1.23}{px} misalignment. The reason is a modality gap: the \emph{local} cross-modal
gradient normalised cross-correlation between SAR and optical is near zero
(NCC~${\approx}0.02$). The misalignment is real and \emph{globally} detectable (rigid phase
correlation), but it carries no \emph{dense, local} similarity signal for a learned field to
climb. A direct local-NCC registration objective likewise produced gradient but could not
reduce its loss. \textbf{Conclusion: dense unsupervised SAR$\to$optical registration cannot be
learned from image similarity; misalignment must be handled by a shift-tolerant supervision
signal rather than by correction.}

\subsubsection{Negative result 2: distributional losses rank sharpness but do not synthesise it}
\label{sec:hfd:neg2}

If misregistration cannot be corrected, the alternative is to supervise detail with a
\emph{shift-invariant} loss. We tested a sliced-Wasserstein (SWD) loss on the Haar detail
sub-bands (matching the sorted coefficient distribution, hence permutation- and
shift-invariant). At the \emph{loss level} the property held decisively: under a synthetic
\SI{1.5}{px} shift, $L_1$ on detail \emph{prefers a blurred prediction} over a
sharp-but-shifted one ($0.037 < 0.056$), whereas SWD \emph{prefers the sharp prediction}
($0.0095 < 0.025$, a $2.6\times$ margin) and is shift-invariant (shift-sensitivity index
$0.13$). This mechanistically explains the blur: pixel losses reward smearing when the target
is misaligned.

However, a fair held-out generalisation A/B (Section~\ref{sec:hfd:exp}) refuted the loss as a
\emph{training driver}: replacing detail $L_1$ with detail SWD made the output marginally
\emph{softer}, not sharper (mid-freq power $0.629$ vs.\ $0.655$; detail-band power $0.582$ vs.\
$0.626$; LPIPS $0.311$ vs.\ $0.305$). \textbf{Conclusion: a loss that \emph{ranks} a sharp
output above a blurry one does not, by minimisation, \emph{produce} sharp output.} Matching a
coefficient histogram fixes the energy budget but supplies no spatial gradient to place
coherent structure; the perceptual and adversarial terms then dominate. This points to the
operative principle behind the final method: only an \emph{adversarial} signal (or a
position-aligned pixel signal, which is unavailable here) carves coherent high-frequency
structure.

\subsection{Method: a high-frequency discriminator}
\label{sec:hfd:hfd}

We introduce a dedicated \emph{high-frequency discriminator} (HF-D) that supplies adversarial
pressure specifically on the deficient frequency band. Let $h(x) = x - g_\sigma * x$ be the
high-pass residual of an image, where $g_\sigma$ is a Gaussian kernel ($\sigma=2$). The HF-D,
$D_H$, is a conditional spectral-normalised PatchGAN that receives the SAR condition $s$
concatenated with the high-pass of the optical image:
\begin{equation}
D_H\big([\,s,\; h(y)\,]\big),
\end{equation}
and is trained adversarially against the generator output $\hat{y}=G(s)$. With LSGAN losses,
the discriminator and generator high-frequency terms are
\begin{align}
\mathcal{L}_{D_H} &= \tfrac{1}{2}\,\mathbb{E}\big[(D_H([s,h(y)])-1)^2\big]
                   + \tfrac{1}{2}\,\mathbb{E}\big[D_H([s,h(\hat{y})])^2\big],\\
\mathcal{L}^{G}_{H} &= \mathbb{E}\big[(D_H([s,h(\hat{y})])-1)^2\big],
\end{align}
and the generator objective adds $\lambda_H\,\mathcal{L}^{G}_{H}$ to the existing loss stack
(adversarial main discriminator, feature matching, MS-SSIM, LPIPS, focal-frequency, and
per-band $L_1$). Because the HF-D judges the \emph{realism of detail} rather than its
position-by-position match, it is tolerant to the residual misregistration that defeats pixel
supervision, while still imposing the coherence that distributional losses lack. Operating on
the high-pass residual focuses the discriminator's capacity on exactly the mid/high-frequency
octaves identified as deficient in Table~\ref{tab:errorprofile}.

Spectral normalisation on every convolution is essential: an earlier sub-band discriminator
without it diverged (logits ${\sim}10^5$). With spectral normalisation the HF-D is stable
(Section~\ref{sec:hfd:exp}). The HF-D is additive to, and independent of, the main conditional
discriminator; setting $\lambda_H=0$ recovers the baseline exactly, giving a clean single-knob
ablation.

\subsection{Experimental verification}
\label{sec:hfd:exp}

\paragraph{Protocol.} Because the target effect is a \emph{generalisation} phenomenon
(misalignment-induced blur appears only when the model cannot memorise each pair's specific
shift), single-batch overfitting is the wrong probe---it favours pixel $L_1$ by memorisation.
We therefore evaluate on a \emph{held-out} A/B: from a common warm-started initialisation, two
arms are trained for \num{2200} generator steps on raw (unregistered) SEN1-2 data, identical in
every respect except the variable under test, and measured on a fixed held-out set of
\num{32} pairs never seen in training. Sharpness is read from the radial mid-frequency power
ratio (target $1.0$); coherence from the high-frequency power ratio (values $\gg 1.15$ indicate
incoherent noise); perceptual quality from LPIPS. PSNR/SSIM are reported but are \emph{not}
used as sharpness judges, since under misaligned ground truth they reward blur.

\paragraph{Results.} Table~\ref{tab:hfd_ab} reports the held-out A/B. The HF-D arm increases the
mid-frequency power ratio at every checkpoint and raises detail-band energy toward the target,
while the high-frequency ratio rises but remains below $1.0$---the added detail is coherent, not
noise. LPIPS is unchanged and the discriminator is stable throughout ($d$-loss
$0.23\!\rightarrow\!0.25$, no divergence). The same harness refuted the distributional (SWD)
loss, which moved the metrics the wrong way.

\begin{table}[t]
\centering
\caption{Held-out generalisation A/B after \num{2200} steps on raw (misaligned) SEN1-2.
Mid-freq.\ and detail power ratios target $1.0$ (higher $=$ sharper); high-freq.\ ratio should
approach $1.0$ from below (values $\gg 1.15$ indicate incoherent noise). Baseline $=$ full loss
stack; HF-D $=$ baseline $+\,\lambda_H\mathcal{L}^G_H$. The rejected sliced-Wasserstein arm is
shown for contrast.}
\label{tab:hfd_ab}
\begin{tabular}{lccc}
\toprule
Metric (held-out) & Baseline & \textbf{HF-D} & SWD (rejected) \\
\midrule
Mid-freq.\ power ratio $\uparrow$   & $0.652$ & $\mathbf{0.695}$ & $0.629$ \\
Detail-band power ratio $\uparrow$  & $0.626$ & $\mathbf{0.687}$ & $0.582$ \\
High-freq.\ ratio ($\to 1.0$)       & $0.643$ & $0.696$          & --- \\
LPIPS $\downarrow$                  & $0.307$ & $0.307$          & $0.311$ \\
PSNR (not a sharpness judge)        & $17.08$ & $16.98$          & $17.14$ \\
SSIM                                & $0.298$ & $0.294$          & $0.300$ \\
\bottomrule
\end{tabular}
\end{table}

\paragraph{Scope of the claim.} The \num{2200}-step A/B establishes the \emph{sign} and
\emph{coherence} of the effect (HF-D sharpens, the added energy is coherent, no perceptual or
stability regression), not its saturated magnitude; the absolute margins at this horizon are
modest (mid-freq.\ ${+}0.04$) because both arms are far from convergence. The full-length
training results are reported in Table~\ref{tab:hfd_final} (run
\texttt{llwt-v0.5.1-hfd}; \,$\langle$\textsc{fill}$\rangle$).

\begin{table}[t]
\centering
\caption{Final \texttt{llwt-v0.5.1-hfd} versus the $\lambda_H{=}0$ baseline at convergence
($\langle$\textsc{fill after run}$\rangle$).}
\label{tab:hfd_final}
\begin{tabular}{lcccc}
\toprule
Model & PSNR & SSIM & LPIPS & FID \\
\midrule
Baseline ($\lambda_H{=}0$)          & \textsc{fill} & \textsc{fill} & \textsc{fill} & \textsc{fill} \\
\textbf{HF-D} ($\lambda_H{=}1$)     & \textsc{fill} & \textsc{fill} & \textsc{fill} & \textsc{fill} \\
\bottomrule
\end{tabular}
\end{table}

\subsection{Ablation and reproducibility}
\label{sec:hfd:ablation}

The headline ablation is a single flag: \texttt{model.dis.highfreq.enabled} (equivalently
$\lambda_H \in \{0,1\}$), with all other settings, data order, and the warm-start checkpoint
held fixed---an exact single-variable comparison. Secondary controls, all refuted cheaply and
reported as negative results, are: the learned dense aligner (Section~\ref{sec:hfd:neg1}) and
the sliced-Wasserstein detail loss (Section~\ref{sec:hfd:neg2}). All experiments are driven by
config flags in \texttt{src/models/llwt\_v5/config.yaml}; the verification harness
(\texttt{proof\_hfd\_generalize.py}) and the diagnostic scripts are version-controlled.

\subsection{Limitations and future work}
\label{sec:hfd:limits}

The HF-D sharpens by imposing realistic detail statistics, not ground-truth-exact detail; under
misaligned references this is the correct objective, but it means PSNR (which rewards blur here)
is a poor progress metric and may even decrease slightly. The high-frequency power ratio must be
monitored over long training: should it overshoot $1.0$ substantially, $\lambda_H$ should be
reduced (e.g.\ $1.0\!\rightarrow\!0.5$) to avoid re-introducing incoherent texture. The
one-to-many ambiguity of low-texture (water) regions is mitigated but not solved by adversarial
pressure alone; an explicit aleatoric-uncertainty head, or semantic conditioning from a frozen
self-supervised optical encoder, is a natural extension. Finally, the modality-gap result
(Section~\ref{sec:hfd:neg1}) suggests that any future alignment component should operate on a
global/low-degree-of-freedom transform supervised by an offline rigid estimate, rather than on a
dense field learned from image similarity.
